\documentclass[11pt]{article}
\usepackage[margin=1in]{geometry}
\usepackage{amsmath,amssymb,amsthm}
\usepackage{booktabs}
\usepackage{hyperref}
\usepackage{xcolor}
\usepackage{algorithm}
\usepackage{algpseudocode}
\usepackage{microtype}

\hypersetup{colorlinks=true, linkcolor=blue, citecolor=blue, urlcolor=blue}

\newcommand{\pT}{p_{\mathrm{T}}}
\newcommand{\pS}{p_{\mathrm{S}}}
\newcommand{\KL}{\mathrm{KL}}
\newcommand{\JSD}{\mathrm{JSD}}
\newcommand{\logpT}{\log p_{\mathrm{T}}}
\newcommand{\logpS}{\log p_{\mathrm{S}}}

\title{Selective On-Policy Self-Distillation for Mathematical Reasoning:\\
       Methods and Extensions}
\author{}
\date{}

\begin{document}
\maketitle

\begin{abstract}
On-Policy Self-Distillation (OPSD) trains a language model by distilling knowledge from a
privileged version of itself—one that sees a reference solution at inference time—into a
non-privileged student.  While effective, OPSD applies distillation uniformly across all
token positions and all rollouts, without regard for whether teacher guidance is
\emph{beneficial} at each point.  We describe three extensions that address distinct failure
modes of uniform distillation: (1) \textbf{CH-OPD}, which selects positions where
the teacher's advantage is locally high; (2) \textbf{FED}, which preserves strategy-level
diversity that token-level KL tends to collapse; and (3) \textbf{OG-OPD}
(Outcome-Gated OPD), our primary contribution, which suppresses distillation on
rollouts the student already solved correctly—avoiding corruption of self-discovered solution
paths.  OG-OPD requires no additional generation and is orthogonal to, yet theoretically
complementary with, the SDAR framework~\cite{sdar2025}.
\end{abstract}

% ============================================================
\section{Background: On-Policy Self-Distillation (OPSD)}
\label{sec:opsd}
% ============================================================

\paragraph{Setup.}
Let $\mathcal{M}_\theta$ be a language model with parameters $\theta$, and let $(x, s^*)$
denote a problem–solution pair drawn from a training distribution.  OPSD maintains a
\emph{single} parameter vector $\theta$ but creates two \emph{input contexts} for the same
forward pass:
\begin{align*}
  c_S &= \texttt{[problem: } x\texttt{]}
       \quad\text{(student context)},\\
  c_T &= \texttt{[problem: } x\texttt{, reference solution: } s^*\texttt{]}
       \quad\text{(teacher context)}.
\end{align*}
At each training step, the student samples a completion $y = (y_1,\dots,y_T)$ from
$\mathcal{M}_\theta(\cdot \mid c_S)$.  Both the student and the teacher then score this
\emph{same} completion:
\[
  \pS_t = \mathcal{M}_\theta(\cdot \mid c_S, y_{<t}), \qquad
  \pT_t = \mathcal{M}_\theta(\cdot \mid c_T, y_{<t}).
\]
The student and teacher share weights $\theta$; only their conditioning contexts differ.

\paragraph{Loss.}
OPSD minimizes the token-level generalized Jensen--Shannon divergence:
\begin{equation}
  \mathcal{L}_{\mathrm{OPSD}}(\theta)
  = \frac{1}{|\mathcal{T}|}\sum_{t \in \mathcal{T}}
    D_{\JSD}^{(\beta)}\!\left(\pT_t \,\|\, \pS_t\right),
  \label{eq:opsd}
\end{equation}
where $\mathcal{T}$ is the set of completion token positions (prompt positions are masked
with label $-100$), and the $\beta$-interpolated JSD is
\begin{equation}
  D_{\JSD}^{(\beta)}(p \| q)
  = \beta\,\KL(p \| m) + (1-\beta)\,\KL(q \| m),
  \qquad m = \beta p + (1-\beta)q.
\end{equation}
Setting $\beta=0$ recovers reverse KL, $\beta=1$ forward KL, and $\beta=0.5$ the symmetric
JSD.  The paper uses $\beta=0$ (reverse KL, i.e.\ the student mimics the teacher's
probability mass).

\paragraph{On-policy generation.}
Because $y$ is sampled from the \emph{current} student at each step, the training
distribution $\mathcal{M}_\theta(\cdot \mid c_S)$ always matches the model's current
behavior.  This avoids the distribution-shift problem of offline distillation, where a fixed
set of teacher completions becomes stale as $\theta$ evolves.

\paragraph{Rollout scoring.}
Each rollout $y^{(i)}$ is also passed to a symbolic verifier $\mathcal{V}(y^{(i)}, s^*)
\in \{0,1\}$, producing a binary reward $r_i$.  In the base OPSD formulation the reward is
tracked for monitoring but \emph{not} used in the distillation loss.

% ============================================================
\section{Gap 1: Uniform Distillation Ignores Where the Hint Helps}
\label{sec:ch-opd}
% ============================================================

OPSD applies Eq.~\eqref{eq:opsd} at every token in every rollout.  However, the reference
solution $s^*$ changes the teacher's prediction only at a subset of positions—those where
knowing the answer causally alters the next-token distribution.  At positions where
$\pT_t \approx \pS_t$ (the student would have predicted similarly anyway), the JSD loss is
already near zero, so distilling there adds training noise without useful signal.

\subsection{Causal-Hinge OPD (CH-OPD)}

CH-OPD~\cite{ch-opd} introduces a per-position \emph{benefit signal}:
\begin{equation}
  B_t = \sum_{a \in \mathcal{V}_t}
        \bigl[\pT_t(a) - \pS_t(a)\bigr] \cdot \hat{V}(a;\, y_{<t},\, s^*),
  \label{eq:benefit}
\end{equation}
where $\mathcal{V}_t$ is a candidate action set (union of the top-$k$ tokens under $\pT_t$
and $\pS_t$, plus the sampled token), and $\hat{V}(a;\cdot)$ is an estimate of the expected
verifier reward when token $a$ is forced at position $t$.  Distillation is applied only at
positions where $B_t > \tau$:
\begin{equation}
  \mathcal{L}_{\mathrm{CH-OPD}}
  = \frac{1}{|\mathcal{H}|}\sum_{t\,:\,B_t>\tau}
    D_{\JSD}^{(\beta)}\!\left(\pT_t \,\|\, \pS_t\right),
  \qquad \mathcal{H} = \{t : B_t > \tau\}.
\end{equation}

\paragraph{Full benefit estimation (expensive).}
$\hat{V}(a;\cdot)$ is estimated by branch-probe sampling: for each candidate token $a$,
prepend $y_{<t} \circ a$ as a forced prefix, generate $n_{\text{probes}}$ short
continuations, and average the verifier rewards.  This requires $|\mathcal{V}_t| \times
n_{\text{probes}}$ extra forward passes per selected position, adding 10--20$\times$
generation overhead per training step.

\paragraph{Logit-divergence proxy (efficient, but a near-no-op).}
A natural zero-cost proxy replaces $\hat{V}$ with the token-level KL divergence:
\begin{equation}
  B_t^{\mathrm{KL}} = \KL\!\left(\pT_t \,\|\, \pS_t\right).
  \label{eq:kl-proxy}
\end{equation}
While this identifies positions where the teacher and student disagree most, it does
\emph{not} capture whether the teacher's preferred tokens lead to better outcomes.
Moreover, it is largely \emph{equivalent to standard OPSD}: positions with small
$\KL(\pT_t\|\pS_t)$ already contribute a small loss in Eq.~\eqref{eq:opsd}, so masking
them with $B_t^{\mathrm{KL}} \leq \tau$ changes the gradient only marginally.  In practice
(with default $\tau=0$) the hinge rate is 1.0—all tokens pass—and the method is
indistinguishable from OPSD.  The genuine novelty of CH-OPD therefore resides entirely in
the full branch-probe mode.

\paragraph{Relationship to SDAR.}
SDAR~\cite{sdar2025} applies a soft token-level gate:
\begin{equation}
  g_t = \sigma\!\left(\beta_{\mathrm{SDAR}} \cdot \Delta_t\right),
  \qquad
  \Delta_t = \logpT(y_t \mid c_T) - \log p_\theta(y_t \mid c_S),
  \label{eq:sdar-gate}
\end{equation}
where $\Delta_t$ is the teacher-minus-student log-probability gap \emph{at the generated
token} $y_t$.  Unlike Eq.~\eqref{eq:kl-proxy}, the SDAR gate is signed: when
$\Delta_t < 0$ (the student assigned higher probability to $y_t$ than the teacher did),
$g_t < 0.5$ and distillation is attenuated.  SDAR uses RL (GRPO) as its primary
objective and the gated distillation as a small auxiliary term ($\lambda=0.01$); it does not
use trajectory-level reward to gate the distillation.

% ============================================================
\section{Gap 2: Token-Level KL Collapses Strategy Diversity}
\label{sec:fed}
% ============================================================

OPSD's token-level JSD treats all reasoning strategies as interchangeable: it pushes every
completion token toward the teacher's distribution, regardless of which high-level solution
approach the completion follows.  For mathematical reasoning, a single problem often admits
multiple valid strategies (algebraic manipulation, geometric argument, number-theoretic
insight, combinatorial counting).  A good policy should maintain probability mass across all
valid strategies, maximizing pass@$k$.

In practice, when the reference solution $s^*$ reflects one particular strategy, OPSD's
token-level loss progressively concentrates the student's output distribution on that
strategy's token patterns, collapsing diversity and reducing pass@$k$.

\subsection{Functional-Equivalence Distillation (FED)}

FED operates at the level of \emph{strategy classes} rather than individual tokens.

\paragraph{Strategy classification.}
Each rollout $y^{(i)}$ is assigned to a strategy class $c^{(i)} \in \mathcal{C}$
(e.g., algebra, geometry, number theory, combinatorics, sequences, trigonometry) by
keyword matching on the solution text.  The verifier reward $r_i$ provides the class-level
quality signal.

\paragraph{Target distribution.}
Let $Q(c) \propto \exp(\rho \cdot \bar{r}_c)$ be a Boltzmann distribution over strategy
classes, where $\bar{r}_c$ is the empirical mean reward within class $c$ and $\rho > 0$ is
a temperature.  The class-level distillation loss pushes the student's strategy distribution
toward $Q$:
\[
  \mathcal{L}_{\mathrm{class}} = \KL\bigl(Q \,\|\, P_\theta\bigr),
\]
where $P_\theta(c)$ is estimated from the rollout frequencies and logits.

\paragraph{Within-class diversity.}
To prevent collapse \emph{within} a high-reward class, FED adds a regularizer against a
frozen reference policy $p_{\mathrm{ref}}$ (the model's initial weights):
\[
  \mathcal{L}_{\mathrm{within}}
  = \lambda_{\mathrm{within}}
    \sum_{c\,:\,\bar{r}_c \geq \tau_{\mathrm{value}}}
    \KL\!\left(p_\theta(\cdot \mid c_T,\, y^{(i)}) \;\|\; p_{\mathrm{ref}}(\cdot \mid c_S,\, y^{(i)})\right).
\]

\paragraph{Efficient strategy discovery.}
The full mode generates $n_{\mathrm{cont}}$ additional continuations per anchor position
(4--5$\times$ extra generation).  The efficient mode reuses the $n_{\mathrm{rollouts}}$
completions already sampled during the student rollout phase: their stochastic decoding
naturally produces distinct strategies, their rewards are already computed, and no extra
generation is needed.

% ============================================================
\section{Gap 3: Distilling Successful Rollouts Can Corrupt Self-Discovered Paths}
\label{sec:og-opd}
% ============================================================

Both OPSD and the methods above share a common assumption: \emph{all rollouts should be
distilled toward the teacher}.  We identify a subtle but important failure mode of this
assumption.

\paragraph{The problem.}
The teacher's privileged context $c_T$ includes $s^*$, meaning the teacher is anchored to a
\emph{specific} reference solution style.  When the student independently discovers a
correct solution (verifier returns $r_i = 1$), that solution may follow a different
reasoning path or stylistic convention than $s^*$.  Applying OPSD's JSD loss on this
rollout pushes the student's token distributions toward the teacher's style—even though the
student's path was \emph{already correct}.  This can:
\begin{enumerate}
\item Erase a solution strategy the student self-discovered, reducing diversity.
\item Introduce gradient noise when the teacher's and student's correct solutions differ in
      surface form but not in mathematical validity.
\item Slow convergence by fighting the RL-like signal implicit in correct rollouts.
\end{enumerate}

\paragraph{The original CH-OPD connection.}
CH-OPD's benefit signal $B_t$ (Eq.~\ref{eq:benefit}) was designed to detect exactly this:
when $B_t < 0$, the teacher's preferred tokens lead to \emph{worse} expected outcomes than
the student's tokens, and distillation is harmful.  However, estimating $B_t$ requires
expensive branch probing, and the logit-KL proxy misses the value component entirely.  We
propose a trajectory-level approximation that costs nothing.

\subsection{Outcome-Gated OPD (OG-OPD)}

\paragraph{Formulation.}
Let $\{(y^{(i)}, r_i)\}_{i=1}^{N}$ be the $N = n_{\mathrm{batch}} \times n_{\mathrm{rollouts}}$
rollouts generated at a training step, with binary rewards $r_i \in \{0,1\}$.
OG-OPD reweights the per-rollout distillation loss by the \emph{failure indicator}:
\begin{equation}
  \boxed{
  \mathcal{L}_{\mathrm{OG-OPD}}
  = \frac{1}{\sum_i (1-r_i)}
    \sum_{i=1}^{N} (1-r_i)\cdot
    \frac{1}{|\mathcal{T}_i|}\sum_{t \in \mathcal{T}_i}
    D_{\JSD}^{(\beta)}\!\left(\pT_t^{(i)} \,\|\, \pS_t^{(i)}\right).
  }
  \label{eq:og-opd}
\end{equation}
In words: rollouts with $r_i=1$ (student solved the problem) receive zero distillation
weight; rollouts with $r_i=0$ (student failed) receive full weight.

\paragraph{Intuition.}
The gate $(1-r_i)$ is a trajectory-level proxy for the CH-OPD benefit: if the student
succeeded, its deviation from the teacher was either irrelevant or \emph{beneficial}, so we
should not push it back.  If the student failed, teacher guidance is likely to help.

\paragraph{Emergent curriculum.}
Early in training, when the model cannot yet solve problems, all $r_i = 0$ and
Eq.~\eqref{eq:og-opd} reduces to standard OPSD—maximum teacher guidance.  As the model
improves and some rollouts succeed, distillation concentrates on the remaining failures,
providing a natural curriculum without any scheduling hyperparameter.

\paragraph{Implementation.}
OG-OPD requires no architectural changes and no extra forward passes beyond standard OPSD.
The rewards $\{r_i\}$ are computed by the verifier during the rollout phase, which is
already part of the OPSD training loop.  The only code change is multiplying each rollout's
loss contribution by $(1-r_i)$ before averaging.

\paragraph{Edge case.}
If all rollouts succeed ($\sum_i (1-r_i) = 0$), the distillation loss is undefined.
In practice, we fall back to standard OPSD in this case (i.e., we set all weights to 1),
treating a step where the student solves every problem as a signal to maintain the
teacher's distributional guidance:
\begin{equation}
  w_i = \begin{cases}1 - r_i & \text{if } \sum_j (1-r_j) > 0, \\ 1 & \text{otherwise.}\end{cases}
\end{equation}

% ============================================================
\section{Comparison of Methods}
\label{sec:comparison}
% ============================================================

Table~\ref{tab:comparison} summarizes the four methods along the key axes that distinguish
them.

\begin{table}[h]
\centering
\caption{Comparison of OPSD variants.}
\label{tab:comparison}
\renewcommand{\arraystretch}{1.3}
\begin{tabular}{@{}lcccc@{}}
\toprule
\textbf{Property} & \textbf{OPSD} & \textbf{CH-OPD} & \textbf{SDAR} & \textbf{OG-OPD (ours)} \\
\midrule
Gating level & None & Per-token & Per-token & Per-trajectory \\
Gating signal & — & $B_t$ (value-weighted) & $\Delta_t$ (log-prob gap) & $r_i$ (verifier reward) \\
Uses trajectory outcome? & No & Indirectly (via $\hat{V}$) & No & Yes, directly \\
Extra generation & None & $10$–$20\times$ & None & None \\
RL objective & No & No & Yes (GRPO primary) & No \\
Handles wrong teacher? & No & Yes (via $B_t$) & Yes (soft, via $g_t$) & Partially (trajectory-level) \\
Protects successful student paths? & No & Yes (via $B_t < 0$) & No & Yes \\
\bottomrule
\end{tabular}
\end{table}

\paragraph{SDAR vs.\ OG-OPD.}
The two methods are \emph{orthogonal}: SDAR asks ``at this token, does the teacher
locally outperform the student?''\ whereas OG-OPD asks ``for this entire rollout, did the
student succeed?''.  Their gating conditions do not imply each other:
\begin{itemize}
  \item A token may have $\Delta_t > 0$ (teacher preferred it locally) even in a
        \emph{successful} rollout that OG-OPD would suppress.
  \item A failed rollout that OG-OPD distills fully may contain tokens with $\Delta_t < 0$
        (student was locally more confident) that SDAR would down-weight.
\end{itemize}
A combined method that gates on \emph{both} conditions—$(1-r_i) \cdot g_t$—would address
both failure modes simultaneously, at no extra generation cost (the $g_t$ computation only
requires the logits already produced in the OPSD forward passes).

% ============================================================
\section{Experimental Setup}
\label{sec:experiments}
% ============================================================

\paragraph{Model and data.}
All experiments use \texttt{Qwen3-1.7B} as the base model.  Training data is the OPSD
paper's dataset \texttt{siyanzhao/Openthoughts\_math\_30k\_opsd} ($\approx$29{,}000
math olympiad problems with verified solutions).  We train on 2{,}000 randomly sampled
problems for 300 steps to study training dynamics.

\paragraph{Hyperparameters.}
Batch size 1 per GPU, 4 rollouts per problem, 4-GPU DDP, gradient accumulation steps 4
(effective batch: 16 rollouts/step).  Learning rate $5 \times 10^{-6}$, AdamW with weight
decay 0.01, gradient clipping at 1.0.  Maximum completion length 1024 tokens.
Gradient checkpointing enabled to reduce activation memory.

\paragraph{Evaluation.}
We evaluate on AIME 2024 (30 problems) using the paper's protocol: temperature 1.0,
top\_p 0.8, max\_tokens 32{,}768, 12 samples per problem, vLLM with tensor parallelism 4.
The primary metric is pass@1 (average fraction correct across samples).

\paragraph{Baselines.}
\begin{itemize}
  \item \textbf{OPSD}: base method (Eq.~\ref{eq:opsd}).
  \item \textbf{OG-OPD}: proposed method (Eq.~\ref{eq:og-opd}).
  \item \textbf{Qwen3-1.7B (no training)}: baseline from the OPSD paper (11.9\% AIME24
        avg@12).
\end{itemize}

\bibliographystyle{plain}
\begin{thebibliography}{9}

\bibitem{opsd}
OPSD Authors.
\newblock On-Policy Self-Distillation for Mathematical Reasoning.
\newblock \emph{arXiv preprint}, 2025.

\bibitem{ch-opd}
Internal project.
\newblock Causal-Hinge OPD: Benefit-Selective Token Distillation.
\newblock 2025.

\bibitem{sdar}
SDAR Authors.
\newblock SDAR: Self-Distilled Agentic Reinforcement Learning.
\newblock \emph{arXiv:2605.15155}, 2025.

\bibitem{selectkd}
Huang et al.
\newblock SelecTKD: Token-Selective Knowledge Distillation via Teacher Entropy.
\newblock 2025.

\end{thebibliography}

\end{document}
